% Source: https://artofproblemsolving.com/community/c3h2730579_similar_triangles_in_some_order
% Stable identity: aops-topic-2730579-q01
First, the correspondence $O\to A$, $A\to C$, $C\to B$ suggested by the names is impossible. It would give $\angle AOC=\angle CAB$, and with $\angle OAC=180^\circ-\angle CAB$ the angles of triangle $OAC$ at $O$ and $A$ would already add up to $180^\circ$.

Similar triangles $OAC$ and $ACB$ share the vertex $A$, whose angles $\angle OAC$ and $\angle CAB$ are supplementary. Checking the possible correspondences (by solving the length equations; all others lead to $C$ on the line or to contradictions), the only one that works is $\triangle OAC\sim\triangle CAB$ ($O\to C$, $A\to A$, $C\to B$). It requires $\angle OAC=\angle CAB$; being supplementary, both are $90^\circ$. Then $\dfrac{OA}{CA}=\dfrac{AC}{AB}$, so $AC^2=OA\cdot AB=54\cdot24=1296$ and $AC=36$.

\begin{asy}
usepackage("amsmath");
size(5cm);
defaultpen(fontsize(10pt));
pair O=(0,0),A=(5.4,0),B=(7.8,0),C=(5.4,3.6);
draw(O--B,linewidth(0.8));draw(O--C--B,linewidth(0.8));draw(A--C,dashed);
draw(A+(0.25,0)--A+(0.25,0.25)--A+(0,0.25));
label("$O$",O,SW);label("$A$",A,S);label("$B$",B,SE);label("$C$",C,N);
label("$54$",(O+A)/2,S,fontsize(8pt));label("$24$",(A+B)/2,S,fontsize(8pt));label("$36$",(A+C)/2,W,fontsize(8pt));
\end{asy}

Triangle $OCB$ has base $OB=78$ and height $AC=36$, so its area is $\dfrac12\cdot78\cdot36=\boxed{1404}$.
